\section{The divisor monoid and its integral bar}
\label{sec:divisor-bar}

The shifts $e_n\mapsto e_{pn}$ compose according to multiplication of
integers.  Fix a finite ordered set $F$ of primes and let
\[
 M_F=\left\{\prod_{p\in F}p^{r_p}:r_p\in\mathbb N_0\right\}.
\]
Unique factorization identifies this monoid with the effective divisors
$\sum_{p\in F}r_p[p]$.  Over any commutative ring $R$, its monoid algebra is
\[
 A_F=R[M_F]=R[u_p:p\in F],\qquad
 \epsilon(u_p)=0,\qquad I_F=\ker\epsilon.
\]
Here $u_p$ represents the divisor $[p]$ and multiplication adds divisors.
The augmentation takes the coefficient of the empty divisor.
The unit is the empty divisor. The coefficient modules have the discrete
topology, and all tensors in this section are algebraic.  All algebra
elements have homological degree zero.

For one prime, multiplication by $u_p$ gives an exact resolution
\[
 0\longrightarrow A_{\{p\}}\eta_p
 \xrightarrow{\ d\ }A_{\{p\}}
 \xrightarrow{\epsilon}R\longrightarrow0,
 \qquad d\eta_p=u_p.
\]
It is exact over arbitrary $R$: multiplication by an indeterminate is
injective, and the quotient is $R$.  The classical Koszul resolution for these polynomial generators is their
finite tensor product:
\[
 K_F=A_F\otimes_R\Lambda_R(\eta_p:p\in F),\qquad |\eta_p|=1,
 \qquad d\eta_p=u_p,
\]
with the graded derivation differential.  The exterior algebra imposes
$\eta_p^2=0$ also in characteristic two.

The normalized bar computing $R\otimes_{A_F}^{\mathbf L}R$ is the
algebraic direct sum
\[
 B(A_F)=\bigoplus_{n\geq0}I_F^{\otimes_R n},
 \qquad |[a_1|\cdots|a_n]|=n,
\]
with
\[
 d[a_1|\cdots|a_n]
 =\sum_{j=1}^{n-1}(-1)^j
       [a_1|\cdots|a_ja_{j+1}|\cdots|a_n].
\]
The empty word has degree zero.  Deconcatenation is its coproduct.

\begin{proposition}[Integral prime-set bar comparison]
\label{prop:divisor-bar}
Give $\Lambda_R(\eta_p:p\in F)$ zero differential and primitive generators
for its graded coalgebra structure.  The map
\[
 \eta_{p_1}\wedge\cdots\wedge\eta_{p_n}
 \longmapsto
 \sum_{\sigma\in S_n}\operatorname{sgn}(\sigma)
       [u_{p_{\sigma(1)}}|\cdots|u_{p_{\sigma(n)}}]
 \quad(p_1<\cdots<p_n)
\]
is a quasi-isomorphism of coaugmented differential graded coalgebras
\[
 \Lambda_R(\eta_p:p\in F)\longrightarrow B(A_F).
\]
It is compatible with inclusions of finite prime sets and with scalar
extension.  In particular,
\[
 \operatorname{Tor}^{A_F}_n(R,R)=\Lambda_R^n(R^F).
\]
\end{proposition}

\begin{proof}
The tensor product defining $K_F$ resolves $R$.  Indeed, each one-variable
resolution is a complex of free $R$-modules with an $R$-linear contraction
onto its constant term.  Their finite tensor product has the same property.
Its terms are free $A_F$-modules.

Lift the displayed map to $K_F\to A_F\otimes_R B(A_F)$ by retaining the
left coefficient.  The target is the normalized left bar resolution of
$R$ over $A_F$.  Its differential has the additional first face
$a[a_1|\cdots|a_n]\mapsto aa_1[a_2|\cdots|a_n]$.
The internal faces in the alternating sum cancel in pairs, because
$u_pu_q=u_qu_p$.  The first faces are exactly the terms of the Koszul
differential.  Thus this lift is a chain map inducing the identity on $R$.
Both complexes are free resolutions: the bar resolution has the usual
contraction inserting the augmentation-ideal part of the left coefficient.
A chain map lifting the identity between projective resolutions is a
chain homotopy equivalence.  To construct its inverse and the two
homotopies, lift each cycle through the next surjective differential,
inductively in homological degree.  Projectivity permits every lift.
Tensoring over $A_F$ with $R$ preserves this homotopy equivalence and
annuls the Koszul differential.

For the coalgebra assertion, cut each alternating word after its $j$th
entry.  Grouping terms by the subset of generators before the cut gives
exactly the signed subset coproduct of the exterior coalgebra.  This
also verifies the empty cuts and the counit.  The formulas preserve
prime inclusions and coefficients.  No factorial is inverted.
\end{proof}

For two primes the degree-two representative is
$[u_p|u_q]-[u_q|u_p]$.  Its differential is
$-u_pu_q+u_qu_p=0$, and its reduced coproduct is
$[u_p]\otimes[u_q]-[u_q]\otimes[u_p]$.
For three distinct primes the exterior degree-three class is nonzero
whenever $R\ne0$.  The untwisted divisor construction therefore has no
uniform concentration in degrees zero through two.

Filtered union over finite prime sets gives the same comparison for
$A=R[u_p:p\text{ prime}]$ and the direct-sum exterior coalgebra on all
prime generators.  Every bar tensor involves only finitely many monomials
and therefore finitely many primes.  Thus both coalgebras and the comparison
are the filtered colimits of the finite-set constructions.  Filtered
colimits of $R$-modules are exact, so their comparison remains a
quasi-isomorphism.  This uses no product completion in word length.

Over $\mathbb C$, the orthonormal monomial completion of $A_F$ is
$\ell^2(M_F)$.  Multiplication by $u_p$ extends to $S_p$.
The self-adjoint diagonal operator
$N_Fu^{\mathbf r}=(\sum_pr_p\log p)u^{\mathbf r}$, with
$\mathbf r=(r_p)_{p\in F}$, has domain
\[
 \mathcal D(N_F)=\left\{(c_{\mathbf r})\in\ell^2(M_F):
 \sum_{\mathbf r}\left(\sum_{p\in F}r_p\log p\right)^2
                 |c_{\mathbf r}|^2<\infty\right\}.
\]
Its heat trace is
$\prod_{p\in F}(1-p^{-s})^{-1}$ for $\operatorname{Re}s>0$.
This constructs the algebra before its scalar trace.  The polynomial
augmentation has $\epsilon(u_p)=0$, whereas the Gibbs state satisfies
$\varphi_{\beta,p}(S_p^*S_p)=1$ and
$\varphi_{\beta,p}(S_p)=0$.  It is a positive functional, not a
multiplicative augmentation.  Adjoining the Hilbert adjoints therefore
changes the available boundary data for a bar.
